%----------------------------------------------------------------------------------------------------%
%----------------------------------------------------------------------------------------------------%
%%%%%%%%%%%                                              HEADER                                                %%%%%%%%%%%
%----------------------------------------------------------------------------------------------------%
%----------------------------------------------------------------------------------------------------%
%
% Universale header.tex – gilt für Arbeitsblätter, Leistungskontrollen, Handreichungen, Nachschlagewerke.
% Neue Pakete und TikZ-Bibliotheken werden HIER ergänzt (nicht im Dokument) und unten protokolliert.
% Quelle: https://github.com/mq90/latex-skills-schule  (CC BY 4.0, Martin Quandt)
%
% Änderungsprotokoll
% 2026-11-09  zusammengeführt aus allen Handreichungs-Headern: TikZ-Bibliotheken fit, shapes.multipart,
%             arrows.meta, shapes.callouts, circuits.logic.US, circuits.logic.IEC; Pakete array, booktabs,
%             longtable; xskak, chessboard, sudoku, tikzpeople, tipa (safe); MusiXTeX als auskommentierte
%             Option (definiert \cup als Achtelnote um – bricht $A \cup B$)
% 2026-11-11  Naturwissenschaften: circuitikz [european] (Schaltpläne), mhchem [version=4] (\ce),
%             chemfig (Strukturformeln), pgf-PeriodicTable (\pgfPT), genealogytree [all] (Stammbäume).
%             NICHT siunitx: kollidiert mit units (\unit[10]{cm}); siunitx nur in eigenständigen Dokumenten
% 2026-11-11  TikZ-Bibliothek decorations.pathmorphing (Federn, Wellen: decoration=coil/snake)
%






%----------------------------------------------------------------------------------------------------%
%%%%%%%%%%%%%                            Gestaltung des Dokuments                                %%%%%%%%%%
%----------------------------------------------------------------------------------------------------%



%Definition der Dokumentenklasse
\documentclass[12pt]{article} %book, article, report, standalone



%Pakete für eine saubere Kodierung
\usepackage[utf8]{inputenc} %UTF-8 als Zeichenkodierung  ... ABER nicht immer notwendig, je nach Compiler
\usepackage[T1]{fontenc} %Verwendet westeuropäische Schriftkodierung
\usepackage[ngerman]{babel} %Ermöglicht deutsche Übersetzung und Silbentrennung
\usepackage{lmodern} %Verwendet moderne Schriftarten



%Pakete für die Gestaltung des Dokuments
\usepackage{geometry} %Ermöglicht die Ausrichtung des Blattes mit \geometry{...}
\geometry{a4paper, top=15mm, left=15mm, right=20mm, bottom=15mm, headsep=10mm, footskip=10mm} %gewählte Ausrichtung des Blattes
\pagestyle{empty} %Art der Seitennummerierung, z.B.: empty , plain
\usepackage{fancyhdr} %Ermöglicht die Gestaltung einer Kopfzeile für LK's oder KA's 
\usepackage{lscape} %Ermöglicht eventuelles Querformat des Blattes mit \begin{landscape} ... \end{landscape}
\usepackage{pdfpages} %Ermöglicht einbinden PDF-Dokumente bzw. einzelene Seiten mit \includepdf[pages={1,3,5}]{Dateiname} bzw. pages=-
\usepackage[german]{layout} %Ermöglicht es mit den Befehl \layout sich die Default-Layout beschreiben zu lassen (Bild vom Layout)
%----------------------------------------------------------------------------------------------------%








%----------------------------------------------------------------------------------------------------%
%%%%%%%%%%%%                                        Packages                                                  %%%%%%%%%%%
%----------------------------------------------------------------------------------------------------%



%Pakete für die Strukturierung des Dokuments
\usepackage[shortlabels]{enumitem} %Ermöglicht die Nummerierung selbst festzulegen mit \begin{enumerate}[Art der Nummerierung] ...   \end{enumerate}
\parindent=0px %Setzt die Absatzeinrückung auf 0 Pixel und ernfernt damit unnötige Leerzeichen vor dem ersten Wort eines neuen Abschnitts
\usepackage{setspace} %Ermöglicht verschiedene Funktion zur Abstandsregelung, z.B.: \hspace{2.5cm} ; \onehalfspacing
\onehalfspacing %Setzt den Zeilenabstand auf 1,5
\usepackage[style=german]{csquotes}%Ermöglicht den deutschen Zitierstil darstellen zu lassen mt \enquote{Zitat}
\usepackage[breaklinks=true,%Zeilenumbruch der Verlinkung ist bündig
			colorlinks=true,%Links können farbig sein
			linkcolor=blue,%Linkfarbe wird auf blau gesetzt
			urlcolor=darkgray,%Farb der URL wird auf dunkelgrau gesetzt
			plainpages=false,%
			hyperfootnotes=false,%
			pdfpagelabels=false%
		]{hyperref} %Ermöglicht die Verlinkung innerhalb des Dokuments 
\usepackage{multicol} %Ermöglicht gleichgroße Aufteilung des Dokuments in n Spalten mit \begin{multicols}{n} ... \end{multicols}
\usepackage{xcolor,colortbl} %Ermöglicht Zeilen, Spalten und Zellen einer Tabelle zu färben
\usepackage{caption} %Ermöglicht die Bezeichnung und Strukturierung von Abbildungen oder Quellcode
\usepackage{fancybox} % Ermöglicht die Umrandung einzelner Abschnitte
\usepackage[most]{tcolorbox} %Ermöglicht die Nutzung von Colorboxen für Merksätze, Lösungen, Definitionen, usw mit \newtcolorbox
\tcbuselibrary{breakable} %Aufteilen der Colorboxen, welche über mehrere Seiten gehen
\usepackage{ifthen} % Ermöglicht \ifthenelse Befehle
\allowdisplaybreaks %Erlaubt Seitenumbrüche innerhalb von align-Umgebung



%Nützliche Pakete für Symbole und mathematischen Formelsatz
\usepackage{amsmath, amssymb, mathrsfs, amsfonts, amsthm, dsfont} %sämtliche Pakete für mathematische Zeichen, Symbole und Formelsatz
\usepackage{units} %Ermöglicht die Angabe von Einheiten mit \unit[Wert]{Einheit}
\usepackage{gensymb} %Ermöglicht Nutzung von \degree für das macro des Grad-Symbols °
\usepackage[official]{eurosym} %Ermöglicht die Nutzung von \euro{}
\usepackage{pifont} %Ermöglicht die Nutzung verschiedener Fonts, z.B.: \ding{42} , \ding{52} , \ding{56}
\usepackage{tikzsymbols} %Ermöglicht die Nutzung von u.a. Smileys und anderen Symbolen, z.B.: \Winkey , \Smiley[1.5]



%Pakete zur Darstellung von Grafiken
\usepackage{tikz} %\begin{tikzpicture} .... \end{tikzpicture}
\usepackage{pgf, pgfplots} %Ermöglicht 2D und 3D Plof-Funktion
\pgfplotsset{compat=1.17}
\usetikzlibrary{arrows,%Pfeile
			calc,%Rechnen in Koordinaten, z.B. mit ($(1,2)+5*(0,1)$) auf Punkt (1,7) kommen
			angles, plotmarks, quotes, babel, intersections, shapes.geometric, decorations.pathreplacing, decorations.pathmorphing, automata, positioning, shadows, patterns, patterns.meta,%
			fit,%Rahmen um mehrere Knoten (UML, Automaten)
			shapes.multipart,%geteilte Knoten (UML-Klassen)
			arrows.meta,%moderne Pfeilspitzen, z.B. -{Triangle[open]}
			shapes.callouts,%Sprechblasen
			circuits.logic.US, circuits.logic.IEC%Logikgatter (digitale Schaltungen)
			} %nützliche Bibliotheken für TikZ
\usepackage{animate} %animierte Darstellungen
\usepackage{graphicx} %\includegraphics[width=..px,height=..px]{Dateiname}
\usepackage{multirow} %Zellen über mehrere Zeilen mit \multirow{n}{*}{Text}
\usepackage{array} %erweiterte Spaltentypen, z.B. >{\centering\arraybackslash}p{2cm}
\usepackage{booktabs} %schöne Tabellenlinien \toprule, \midrule, \bottomrule
\usepackage{longtable} %Tabellen über mehrere Seiten



%gelegentlich hilfreiche Zusatzpakete
\usepackage{hologo} %\hologo{LaTeX}
\usepackage{suetterl} %Sütterlin als Schriftart mit \mathfrak{Text}
\usepackage{qrcode} %Einbindung eines QR-Codes 
\usepackage{url} %Einbindung eines URL-Links
\usepackage{ccicons} %Creative Commons Symbole
\usepackage{listings} %\begin{lstlisting} ... \end{lstlisting}
\usepackage{struktex} %Struktogrammdarstellung
\usepackage{ulem} %Unterstreichung mit \uline{} ggf. mit \hfill im Inneren für Unterstrich bis zum Ende der Zeile
\usepackage{xskak} %Schach: Partien und Züge mit \newchessgame, \mainline{1. e4 e5} (lädt skak)
\usepackage{chessboard} %Schachbretter mit \chessboard[setfen=...]
\usepackage{sudoku} %Sudoku-Gitter mit \begin{sudoku} |5|3| | ... |. \end{sudoku}
\usepackage{tikzpeople} %Personen-Figuren als TikZ-Knoten, z.B. \node[alice, minimum size=2cm]{};
\usepackage[safe]{tipa} %Lautschrift mit \textipa{...}; [safe], weil tipa sonst \; \: \! \* \s umdefiniert (bricht z.B. \(\;\)) – Langformen siehe latex-lautschrift
\usepackage[european]{circuitikz} %Schaltpläne mit deutschen Symbolen: \begin{circuitikz} \draw (0,0) to[lamp] (2,0); \end{circuitikz}
\usepackage[version=4]{mhchem} %Reaktionsgleichungen mit \ce{2 H2 + O2 -> 2 H2O}
\usepackage{chemfig} %Strukturformeln mit \chemfig{H-C(-[2]H)(-[6]H)-H}, freie Elektronenpaare mit \charge
\usepackage{pgf-PeriodicTable} %Periodensystem mit \pgfPT[...], deutsche Namen mit \pgfPTsetLanguage{de}
\usepackage[all]{genealogytree} %Stammbäume mit \begin{genealogypicture} ... \end{genealogypicture}
%\usepackage{musixtex} %Notensatz – NUR für Musik-Dokumente einkommentieren: definiert u.a. \cup als Note um (bricht $A \cup B$)
%\usepackage[list=true, font=large, labelfont=bf, labelformat=brace, position=top]{subcaption}



%abgeänderte Befehle
\renewcommand{\arraystretch}{1.25} %
%\renewcommand{\familydefault}{\sfdefault} %serifenlose Schriftart
\renewcommand{\lstlistingname}{Quellcode} %Umbenennung
\renewcommand{\lstlistlistingname}{Quellcodeverzeichnis} %Umbenennung
%----------------------------------------------------------------------------------------------------%